\section{Implementation}

\lang is implemented as a small compiler/checker whose main job is to preserve the evidence that a fault-tolerant compilation step would otherwise erase.
The implementation is deliberately organized around three artifacts.
A source program is a list of \texttt{LogicalOp} nodes, which records what the algorithm asks for: preparation, logical gates, error correction, measurement, barriers, and classical branches.
A compiled target program is a tree of \texttt{FTStep} nodes, which records how each request was acquired: direct capability use, a switch path, a resource factory followed by consumption, or a derived region acquisition.
An \texttt{Effect} object travels with every step and gives the conservative account that the checker and the paper tables report.
This separation is the implementation analogue of the calculus judgment: source syntax keeps the user's logical story compact, while target syntax makes capability, resource, certificate, and backend evidence explicit.

The intermediate representation is small enough that the interesting invariants are visible in the data structures.
Each live logical qubit has a \texttt{QType(code, distance, mode)}, so a gate step does not merely name its operands; it also records the code in which those operands are currently protected.
Each target step stores \texttt{input\_codes}, \texttt{output\_codes}, an effect, a certificate level, and a free-form \texttt{details} map that carries rule-specific evidence such as source citations, layout events, produced resource identifiers, or consumed resource identifiers.
For example, a non-native $\gate{T}$ on \code{SurfaceD5} is not compiled into an opaque macro called ``apply T''.
It becomes a parent \texttt{acquire\_gate} step with two auditable children:
\[
  \mathsf{prepare\_resource}\ r:\res{TState};
  \quad
  \mathsf{consume\_resource\_gate}\ r\ \mathsf{for}\ \gate{T}.
\]
The identifier $r$ is part of the sidecar, not an explanatory comment, which is what allows the checker to reject a later attempt to consume the same resource again.

The planner treats acquisition as a local search over typed candidates rather than as a fixed gate-expansion pass.
When it encounters a gate $g$ on qubits currently in code $C$, it first checks whether the rule library contains a direct capability for $(C,g)$.
If not, it constructs candidates of two forms.
A switch candidate changes the operands to a code $C'$ that supports $g$, applies the direct rule there, and switches the operands back to $C$.
A resource candidate produces an appropriate resource state, such as a $\res{TState}$, $\res{BellPair}$, or $\res{CCZState}$, and consumes it for the requested operation.
The default score used by \texttt{best} mode is intentionally simple:
\[
  \mathsf{score}(p) =
  \mathsf{cycles}(p)
  + 0.01\,\mathsf{qubitRounds}(p)
  + 25\,\mathsf{switches}(p)
  + 35\,\mathsf{factories}(p)
  + 10^6\,\epsilon(p)
  + 20\,\mathsf{fail}(p).
\]
The point of this formula is not to claim an optimal cost model.
It gives the prototype a deterministic way to choose among qualitatively different acquisitions while leaving the chosen rule sequence, effect fields, and certificate assumptions available for review.

The hybrid strategy adds one compiler idea that is more useful than a gate-by-gate expansion in the case studies.
When the planner sees a contiguous block of gates and at least two of them need acquisition, it asks whether a single target code supports the whole block.
For the current rule library, \code{Tetra15} supports both $\gate{T}$ and $\gate{CNOT}$, while \code{SurfaceD5} supports neither directly.
The planner can therefore compile the block
\[
  \gate{CNOT};\ \gate{T};\ \gate{CNOT};\ \gate{T}
\]
as one \texttt{acquire\_region}: switch the involved qubits from \code{SurfaceD5} to \code{Tetra15}, run certified transversal gates, and switch them back.
The resulting target tree is useful precisely because it is not flattened into an anonymous cost.
The parent region records the source gate sequence and total effect; the children record each switch, each transversal gate, each layout event, and each certificate level.
Thus a reviewer can see that the transversal $\gate{T}$ and $\gate{CNOT}$ steps are certified through the \code{Tetra15} rules, while the compact \code{SurfaceD5}$\leftrightarrow$\code{Tetra15} bridge used in this prototype remains \Assumed{}.

Effect composition is implemented directly on the same object that is serialized into the sidecar.
Sequential composition adds logical-error upper bounds, cycles, qubit-rounds, decoder latency, gate counts, switch counts, and factory counts; it composes acceptance lower bounds multiplicatively; and it takes the maximum of peak-space fields.
Branch composition takes maxima for upper-bound quantities and minima for lower-bound quantities.
These operations are intentionally conservative, but their direct implementation has an important engineering advantage: the table entries in the paper, the SMT obligations emitted by the formal toolchain, and the checker's step summaries all refer to the same effect fields.
There is no separate presentation-only accounting layer.

The checker is a policy interpreter over the generated sidecar.
It flattens the target tree and checks five classes of obligations.
First, every code mentioned by a step must be supported by the selected backend topology.
Second, layout events must match the backend name, topology, capacity, free-qubit accounting, and switch workspace reservation reported by the layout component.
Third, direct \texttt{apply} steps must have a capability witness in the rule library.
Fourth, switch steps must name a known switch rule and that rule must be supported by the backend topology.
Fifth, resource identifiers must be produced before use and consumed at most once.
The certificate policy is checked at the same time: exploratory mode accepts \Assumed{} steps with warnings, whereas strict mode turns the same warnings into errors.
This policy is the implementation mechanism behind the paper's proof boundary.

The full-stack runner connects the compiler/checker to independent evidence generators.
It emits JSON sidecars for every case-study plan, SMT-LIB files for arithmetic bound checks, protocol certificates for selected GF(2) stabilizer and factory facts, decoder reports, geometry reports, artifact hashes, and a Lean build for the small calculus kernel.
The generated report currently records 53 obligations across code specifications, rule specifications, protocols, decoders, checker results, geometry, and Lean, together with 26 generated artifacts.
The report also carries an explicit certificate-boundary statement: GF(2), Z3, geometry, decoder-table checks, Python validation, and Lean discharge local obligations, while imported physical theorems and prototype assumptions remain visible dependencies.
This is the implementation invariant we rely on in the empirical sections: an accepted exploratory plan is auditable, but it is not silently promoted to a strict physics theorem.
